\section{Triple-Proof Quantum Finality: Summary}

\label{sec:triple-proof}

The triple-proof mechanism is the synthesis of the security model's layers:

\begin{table}[ht]
\centering
\caption{Triple-proof epoch certificate structure.}
\label{tab:triple-proof}
\begin{tabular}{lrrr}
\toprule
\textbf{Proof} & \textbf{Size} & \textbf{Verify Time} & \textbf{Assumption Family} \\
\midrule
BLS aggregate signature & 48\,B & 875\,$\mu$s & Discrete log \\
Groth16-compressed ML-DSA & 192\,B & ${\sim}1$--$3$\,ms (CPU) & Lattice + Hash \\
Ringtail threshold (in Groth16) & 0\,B (embedded) & 0\,$\mu$s (included) & Lattice \\
\midrule
\textbf{Total epoch certificate} & \textbf{248\,B} & \textbf{${\sim}2$--$5$\,ms (CPU)} & \textbf{All three} \\
\bottomrule
\end{tabular}
\end{table}

\begin{itemize}[leftmargin=2em]
  \item Block time: 1\,ms. \textbf{[Tested]} on validator testnet with 100 nodes.
  \item Epoch size: 100 blocks (100\,ms per epoch).
  \item Epoch proof generation: ${\sim}400$\,ms CPU / ${\sim}5$--$15$\,ms GPU (estimated) Groth16 prover~\cite{lux-benchmarks}.
  \item Annual storage for epoch proofs: $248 \times 10 \times 86{,}400 \times 365
        \approx 78$\,MB/year.
\end{itemize}

%% ========================================================================
%%  9. THE INVARIANT
%% ========================================================================
